% Cote 119: offprints, reports, courses, CV and letters (typed and printed part): archive cote 119, site page 290 (PDF page 291).
% typescript. Draft compared: none.
% AI-transcribed by Claude Opus 5.5 (claude-opus-5-5, adaptive thinking, effort high), Message Batches API, batch r119_full (msgbatch_01DwZNm5Y7YsKBWyAN4sFs8w), 2026-11-13; not yet checked by eye.
\scanpage{119/290}
\mspage{- 3 -}

pair \hand{$\geq$} 4 et si $q = \hand{$\langle$} a_1,\dots,a_n \hand{$\rangle$}$ est tel que $w_{\hand{1}}(q) = 0$ (autrement dit,
$a_1\dots a_n$ est un carré), on peut montrer que l'élément $(a_1)\dots(a_{n-1})$ de
$H^{n-1}(k, Z/2Z)$ est un \uline{invariant} de la classe de $q$. Le cas $n = 4$ est
particulièrement intéressant.

\textbf{2.2. Comportement de} $w_1(q)$ \textbf{et} $w_2(q)$ \textbf{par torsion}

Soit $x \hand{$\in$} H^1(k, O(q))$. On associe à $x$ des éléments
\[ \hand{$\delta$}^1(x) \hand{$\in$} H^1(k,Z/2Z) \qquad\text{et}\qquad \hand{$\delta$}^2(x) \hand{$\in$} H^2(k,Z/2Z) \]
de la façon suivante :

$\hand{$\delta$}^1(x)$ est l'image de $x$ dans $H^1(k,Z/2Z)$ par l'application déduite
de l'homomorphisme $\det : O(q) \to \hand{$\{$} \pm 1 \hand{$\}$} = Z/2Z$ ;
$\hand{$\delta$}^2(x)$ est le cobord de $x$ (LN 5, p.I-71) relatif à la suite exacte
de groupes algébriques :
\[ 1 \to Z/2Z \to \overset{\hand{$\sim$}}{O}(q) \to O(q) \to 1 . \]
(Le groupe $\overset{\hand{$\sim$}}{O}(q)$ est un certain revêtement quadratique de $O(q)$ qui
prolonge le revêtement spinoriel $\mathrm{Spin}(q) \to SO(q)$. On peut le caractériser
par la propriété suivante : une symétrie par rapport à un vecteur de
carré $a$ se relève en un élément d'ordre 2 de $\overset{\hand{$\sim$}}{O}(q)$ rationnel sur
le corps $k(\hand{$\sqrt{\phantom{a}}$}a)$.)

Les invariants $\hand{$\delta$}^1(x)$ et $\hand{$\delta$}^2(x)$ permettent de calculer les classes
$w_1$ et $w_2$ de la forme $q_x$ déduite de $q$ par torsion au moyen de $x$ .
On a en effet :

(2.2.1) \quad $w_1(q_x) = w_1(q) + \hand{$\delta$}^1(x)$ \qquad dans $H^1(k,Z/2Z)$ ,
(2.2.2) \quad $w_2(q_x) = w_2(q) + \hand{$\delta$}^1(x).w_1(q) + \hand{$\delta$}^2(x)$ \quad dans $H^2(k, Z/2Z)$.

\textbf{2.3. Les conjectures de Milnor}

Soit $\mathbf{k}^M(k) = \hand{$\bigoplus$} \mathbf{k}^M_n(k)$ \uline{l'anneau de Milnor} (mod 2) de $k$\hand{,} défini au
moyen de symboles multilinéaires $(a_1,\dots,a_n) = (a_1)\dots(a_n)$ , $a_i \hand{$\in$} k^*$, avec
les relations $2(a) = 0$ et $(a,b) = 0$ si $a+b = 1$ .

Soient $W_k$ l'anneau de Witt de $k$, et $I_k$ son idéal d'augmentation\hand{\unsure{,}} noyau de
l'homomorphisme canonique $W_k \to Z/2Z$ .
On définit de façon naturelle des homomorphismes
(2.3.1)\hand{\unsure{.}} \quad $\mathbf{k}^M_n(k) \to I^n_k/I^{n+1}_k$
et
(2.3.2) \quad $\mathbf{k}^M_n(k) \to H^n(k,Z/2Z)$ \hand{\unsure{.}}

\mspage{\hand{290}}
